\documentclass[a4paper]{article}
\usepackage{amsfonts}
\usepackage{amsmath}
\usepackage{amssymb}
\usepackage{graphicx}     

\begin{document}
  
\noindent \large Auteur: Torben Petré \\
\noindent \large Studierichting: Informatica\\
\noindent \large Oefeningenreeks 3A nr. 19.6\\

\medskip

\normalsize

Bereken de afgeleide functie van $f(x) = \dfrac{5}{(x-2)(x+3)}$\\

$= 5(x^2 + x - 6)^{-1}$\\

$= 5 \cdot D((x^2 + x - 6)^{-1}) \cdot (2x + 1)$\\

$= -5 \cdot (x^2 + x - 6)^{-2} \cdot (2x + 1)$\\

$= - \dfrac{10x + 5}{(x^2 + x - 6)^2}$\\

\begin{thebibliography}{999}
%\bibitem{a} Referentie
\end{thebibliography}
\end{document} 
